% ===== PRÉAMBULE CANONIQUE — à coller VERBATIM en tête de CHAQUE .tex =====
\documentclass[11pt,a4paper]{article}
\usepackage[utf8]{inputenc}\usepackage[T1]{fontenc}\usepackage{lmodern}
\usepackage[french]{babel}
\usepackage{amsmath,amssymb,mathtools}\usepackage{icomma}
\usepackage[version=4]{mhchem}          % \ce{...} : équations & formules
\usepackage{siunitx}\sisetup{locale=FR} % \qty{}{}, \si{}, \ang{}
\usepackage{chemfig}                    % \chemfig{...} : structures développées
\usepackage{xcolor}\usepackage{colortbl}
\usepackage{tikz}\usepackage{pgfplots}\pgfplotsset{compat=1.18}
\usepackage[most]{tcolorbox}\usepackage{enumitem}\usepackage{array,booktabs}
\usepackage[a4paper,margin=2.2cm]{geometry}
\usetikzlibrary{arrows.meta,calc,angles,quotes,positioning,patterns,backgrounds,shapes.geometric}
\definecolor{primary}{RGB}{222,49,99}\definecolor{primarydark}{RGB}{150,22,55}
\definecolor{defcol}{RGB}{0,114,178}\definecolor{methcol}{RGB}{52,92,148}
\definecolor{excol}{RGB}{0,135,90}\definecolor{attncol}{RGB}{200,40,55}
\tcbset{boxbase/.style={enhanced,breakable,boxrule=0.9pt,arc=2.5pt,
  left=3mm,right=3mm,top=2.3mm,bottom=2.3mm,fonttitle=\bfseries,coltitle=white}}
% --- boîtes TUTORIEL (noms EXACTS, ne pas renommer) ---
\newtcolorbox{cle}[1][]{boxbase,colback=defcol!6,colframe=defcol,
  title={\textsc{À retenir}\ifx\relax#1\relax\relax\else\textmd{ : \itshape #1}\fi}}
\newtcolorbox{methode}[1][]{boxbase,colback=methcol!7,colframe=methcol,sharp corners,
  title={\textsc{Méthode}\ifx\relax#1\relax\relax\else\textmd{ --- \itshape #1}\fi}}
\newtcolorbox{exemplebox}[1][]{boxbase,colback=excol!5,colframe=excol,
  title={\textsc{Exemple}\ifx\relax#1\relax\relax\else\textmd{ : \itshape #1}\fi}}
\newtcolorbox{piege}[1][]{boxbase,colback=attncol!6,colframe=attncol,
  title={\textsc{Piège}\ifx\relax#1\relax\relax\else\textmd{ : \itshape #1}\fi}}
\newtcolorbox{remarque}[1][]{boxbase,colback=black!4,colframe=black!45,
  title={\textsc{Remarque}\ifx\relax#1\relax\relax\else\textmd{ : \itshape #1}\fi}}
% --- boîtes EXERCICE / CORRIGÉ (noms EXACTS, auto-comptées) ---
\newtcolorbox[auto counter]{exo}[1][]{boxbase,colback=primary!4,colframe=primary,
  title={\textsc{Exercice}~\thetcbcounter\ifx\relax#1\relax\relax\else\textmd{ --- \itshape #1}\fi}}
\newtcolorbox[auto counter]{corrige}[1][]{boxbase,colback=excol!4,colframe=excol,
  title={\textsc{Corrigé}~\thetcbcounter\ifx\relax#1\relax\relax\else\textmd{ --- \itshape #1}\fi}}
\newcommand{\R}{\mathbb{R}}\newcommand{\N}{\mathbb{N}}\newcommand{\Z}{\mathbb{Z}}\newcommand{\ds}{\displaystyle}
% (un \usepackage{fancyhdr}+en-tête est possible ; il est ignoré côté web)
% =========================================================================
\begin{document}
\sisetup{per-mode=symbol}

\begin{corrige}[de la concentration au pH]
$\mathrm{pH} = -\log[\ce{H3O+}]$ ; rappel $\log 10^{-n} = -n$.
\begin{enumerate}[label=\alph*)]
  \item $\mathrm{pH} = -\log 10^{-3} = 3$.
  \item $\mathrm{pH} = -\log 10^{-11} = 11$.
  \item $\mathrm{pH} = -\log 10^{-5} = 5$.
  \item $[\ce{H3O+}] = 0,010 = 10^{-2}$, donc $\mathrm{pH} = 2$.
\end{enumerate}
\end{corrige}

\begin{corrige}[du pH à la concentration]
$[\ce{H3O+}] = 10^{-\mathrm{pH}}$.
\begin{enumerate}[label=\alph*)]
  \item $\qty{e-2}{mol\per L}$.
  \item $\qty{e-9}{mol\per L}$.
  \item $\qty{e-7}{mol\per L}$.
  \item $\qty{e-1}{mol\per L} = \qty{0,1}{mol\per L}$.
\end{enumerate}
\end{corrige}

\begin{corrige}[le produit ionique de l'eau]
On utilise $[\ce{OH-}] = \dfrac{10^{-14}}{[\ce{H3O+}]}$ et $\mathrm{pH}+\mathrm{pOH}=14$.
\begin{enumerate}[label=\alph*)]
  \item $[\ce{OH-}] = \dfrac{10^{-14}}{10^{-4}} = \qty{e-10}{mol\per L}$.
  \item $[\ce{H3O+}] = \dfrac{10^{-14}}{10^{-2}} = \qty{e-12}{mol\per L}$.
  \item $\mathrm{pOH} = 14 - \mathrm{pH} = 14 - 4 = 10$.
\end{enumerate}
\end{corrige}

\begin{corrige}[acide, neutre ou basique ? (à \qty{25}{\celsius})]
Repère par rapport à $\mathrm{pH}=7$ (ou $[\ce{H3O+}]=10^{-7}$).
\begin{enumerate}[label=\alph*)]
  \item $\mathrm{pH}=3 < 7$ : \textbf{acide}.
  \item $\mathrm{pH}=7$ : \textbf{neutre}.
  \item $\mathrm{pH}=11 > 7$ : \textbf{basique}.
  \item $[\ce{H3O+}]=10^{-9} < 10^{-7}$, soit $\mathrm{pH}=9 > 7$ : \textbf{basique}.
\end{enumerate}
\end{corrige}

\begin{corrige}[de $K_a$ au $\mathrm{p}K_a$, et classement des forces]
\begin{enumerate}[label=\alph*)]
  \item $\mathrm{p}K_a = -\log(10^{-5}) = 5$.
  \item Plus le $\mathrm{p}K_a$ est \emph{petit}, plus l'acide est \emph{fort}. Du plus fort au plus faible :
        \[ \ce{HF}\ (3{,}2) \;>\; \ce{CH3COOH}\ (4{,}8) \;>\; \ce{HClO}\ (7{,}5) \;>\; \ce{NH4+}\ (9{,}2). \]
\end{enumerate}
\end{corrige}

\end{document}
